%
% Preface
%

\chapterimage{Ramon_Llull_-_Ars_Magna_Fig_1.pdf} % Chapter heading image

\chapter*{Preface}

\begin{quote}
\begin{flushright}
\emph{Perfection is achieved not when there is nothing more to add, \\
but when there is nothing left to take away.}\\
Antoine de Saint-Exupéry\\
\end{flushright}
\end{quote}
\bigskip

%
% Why: We want to venture into the unknown by transcending the limits of science.
% How: By developing a rigorous mathematical theory based on computability and randomness.
% What: A methodology, and its practical implementation, to minimize humanity's ignorance.
%

% "Breaking AI Limits: Mathematical Methods to Discover the Unknown Unknown"

We live in an era that places a high value on knowledge and celebrates scientific advancement. The development of vaccines in record time, the detection of gravitational waves, and the use of artificial intelligence to model protein structures are just a few examples of the extraordinary progress science has made in recent years. Yet, despite these achievements, our scientific understanding remains incomplete. We continue to face fundamental questions about how knowledge is acquired, how scientific explanations are constructed, and how progress towards a more complete understanding can be measured.

This book is driven by a central aspiration: to overcome the methodological limitations of science in order to challenge established knowledge, explore the unknown, and progress toward the limits of scientific understanding. Our ultimate aim is to investigate both well-known unsolved problems (the known unknowns) and problems that remain outside our current awareness (the unknown unknowns). Achieving this ambition requires a systematic approach to understanding what we do not know, and to determining how scientific inquiry can progressively reduce our ignorance.

Our starting assumption is that ignorance itself can be studied mathematically. To this end, we propose a new mathematical framework, named the Theory of Nescience, grounded in computability and algorithmic information theory. Rather than attempting to measure knowledge directly, the theory approaches scientific understanding through its limitations, seeking to identify and quantify the information that is missing, inadequate, or unnecessary in our current accounts of the phenomena we investigate.

The theory distinguishes three fundamental elements of scientific knowledge: entities, representations, and descriptions. Entities are the objects of scientific investigation, representations encode information about these entities, and descriptions provide computational models of their representations. Within this framework, ignorance arises from two fundamental sources: \emph{miscoding}, which measures the inadequacy of representations, and \emph{mismodel}, which measures the inadequacy of descriptions. Together, these two quantities provide a mathematical characterization of nescience and allow scientific inquiry to be formulated as a multiobjective optimization problem.

A particularly striking consequence of this framework is the relationship between perfect knowledge and randomness. Perfect knowledge requires both an adequate representation of an entity and a minimal description of that representation. No shorter description can reconstruct that representation. In the language of algorithmic information theory, a minimal description is incompressible up to a fixed additive constant, establishing a connection between perfect knowledge and algorithmic randomness. This does not mean that randomness itself constitutes knowledge, but rather that incompressibility is a necessary property of descriptions that have reached their theoretical minimum.

From these foundations emerges a mathematical framework for studying how scientific knowledge is acquired, organized, and refined. The theory allows us to characterize the limitations of existing hypotheses, compare alternative representations and descriptions, investigate the conditions for perfect knowledge, and examine the boundaries separating what is known from what remains unknown.

This first volume is dedicated to establishing the mathematical foundations of the Theory of Nescience. It develops the fundamental concepts, definitions, and theoretical results required to quantify ignorance and analyze the progression of scientific knowledge. The practical application of these ideas to scientific discovery, artificial intelligence, and the systematic investigation of unknown research problems constitutes a broader research program that builds upon these foundations.

Ultimately, the Theory of Nescience invites us to reconsider the aims of scientific inquiry. Instead of striving merely to maximize what we know, it encourages us to minimize what we do not know. By shifting our focus from the accumulation of knowledge to the reduction of ignorance, we open the door to new ways of understanding and measuring scientific progress. This volume establishes the theoretical groundwork for that pursuit.

% Subsection: Research Agenda

\section*{Research Agenda}

The Theory of Nescience is motivated by a collection of fundamental questions concerning the nature of scientific knowledge, the measurement of ignorance, and the limits of scientific understanding. Some of these questions can be addressed through mathematical definitions and theoretical results, while others require further investigation and practical development. This volume concentrates on the mathematical foundations, establishing a framework upon which a broader research program can be developed.

\paragraph{Foundational Questions}

\emph{Can we provide a quantitative characterization of our ignorance regarding a research entity or topic?} Rather than attempting to measure knowledge directly, we investigate whether ignorance can be characterized through deficiencies in the representations and descriptions used in scientific inquiry. A mathematical measure of ignorance would allow us to evaluate scientific hypotheses and investigate how our understanding changes as new knowledge is acquired.

\emph{What are the fundamental sources of scientific ignorance?} Our understanding may be limited because the available representations fail to capture the entities under investigation, because our descriptions fail to account adequately for those representations, or because both are inadequate. Distinguishing these sources would allow us to identify the limitations of scientific hypotheses and determine how they might be improved.

\emph{Can we compare the extent of our ignorance across alternative scientific hypotheses and research areas?} Such comparisons would allow us to evaluate competing explanations and examine differences in scientific understanding. Extending these measurements to research areas could provide a foundation for comparing the state of knowledge across different domains, subject to suitable conditions of comparability.

\emph{Can we characterize scientific progress as the reduction of ignorance?} Scientific inquiry involves the continuous refinement of representations and descriptions. By identifying the conditions under which these refinements reduce ignorance, we may establish mathematical criteria for evaluating scientific progress and selecting among competing hypotheses. Whether these criteria can support universally effective procedures for acquiring knowledge remains a separate question.

\emph{What constitutes perfect knowledge?} Understanding what it means to know an entity completely is essential to defining the limits of scientific understanding. A mathematical characterization of perfect knowledge would establish an ideal reference against which scientific hypotheses could be evaluated. However, defining such a state does not necessarily imply that we can attain it or recognize it in practice.

\emph{Can we characterize the boundaries between what is known and what remains unknown?} Scientific inquiry extends beyond the refinement of established knowledge to the recognition of entities and problems not yet incorporated into scientific understanding. Distinguishing between what is knowable in principle and what has become effectively known may provide a mathematical foundation for investigating both known unknowns and unknown unknowns.

\paragraph{Directions for Future Research}

\emph{Can we devise systematic methods for discovering previously unknown research entities and problems?} Many scientific advances arise from questions that had not previously been formulated. The mathematical characterization of representations and their combinations offers a starting point for generating candidate research entities. Developing reliable procedures to identify, evaluate, and prioritize genuinely novel research questions remains an important objective.

\emph{Can we use quantitative measures of ignorance to improve research prioritization?} Comparing the limitations of existing scientific knowledge could help identify promising directions for investigation. However, effective prioritization would also require measures of scientific relevance, expected benefits, and available resources. Integrating these considerations with nescience could eventually provide a more systematic approach to allocating research effort.

\emph{Can we measure the intrinsic complexity of a research topic and the effort required to understand it?} The amount of ignorance associated with a topic need not reflect its inherent complexity or the resources required to advance our understanding. Developing mathematical measures of these distinct properties could contribute to a more comprehensive characterization of scientific inquiry and its practical limitations.

\emph{Can we establish mathematical criteria for distinguishing science from pseudoscience?} The capacity to formulate hypotheses, evaluate their limitations, and progressively reduce ignorance may provide useful criteria for assessing scientific practices. Whether these properties are sufficient to establish a rigorous distinction between science and pseudoscience remains an open question.

The questions addressed in this volume establish the mathematical foundations of the Theory of Nescience, while the remaining questions outline opportunities for future development. The theory is not intended as a closed doctrine, but as a framework for investigating the nature and limits of scientific knowledge. Its ultimate ambition extends beyond quantifying ignorance: it seeks to provide the foundations for a more systematic understanding of how knowledge is acquired, evaluated, and refined.

% Subsection: Origins of the Theory of Nescience

\section*{Origins of the Theory of Nescience}

It was in 1991, when I was eighteen years old, that I first encountered a statement attributed to Pablo Picasso: \emph{``Computers are useless. They can only give you answers.''} What struck me was not a dismissal of computers, but the distinction between answering questions and identifying questions worth asking. That idea stayed with me for years. In 2014, more than two decades later, I began exploring whether this distinction could be approached from a computational perspective.

My initial goal was to devise a method for identifying interesting and underexplored scientific questions, including those that had not yet been formulated. One particularly restless night, several concepts I had been considering, such as miscoding, relevance, nescience, and the unknown unknowns, suddenly began to fit together. My long-standing interest in information theory and Kolmogorov complexity provided a possible mathematical foundation for these emerging ideas.

Over the following weeks, I conducted a series of computer experiments to investigate these concepts. The results were encouraging, but I gradually realized that discovering new scientific questions required a more fundamental understanding of ignorance itself. I began to wonder whether our lack of knowledge could be quantified, how scientific understanding could be improved, and whether perfect knowledge could be characterized mathematically.

This line of inquiry led to a particularly surprising connection between perfect knowledge and algorithmic randomness. A minimal description is necessarily incompressible up to a fixed additive constant, suggesting that randomness plays an important role in characterizing the theoretical limits of knowledge. Further investigation also revealed the importance of distinguishing between the inadequacy of our representations and the inadequacy of the descriptions used to model them. These ideas eventually became the foundations of the Theory of Nescience.

Over the following years, feedback from colleagues and researchers, together with experiments in data science and machine learning, helped me refine the theory and explore its potential applications. What began as an attempt to discover interesting scientific questions gradually developed into a mathematical framework for investigating the nature of ignorance, the acquisition of knowledge, and the limits of scientific understanding.

This first volume presents the mathematical foundations of that framework. Its broader ambition, however, is to develop methods not only for answering scientific questions, but also for identifying which questions are worth asking, including those that have not yet been imagined.

%
% Section: About the book
%

\section*{About the Book}

The Theory of Nescience draws on concepts from several academic disciplines, particularly discrete mathematics, computability, algorithmic information theory, multiobjective optimization, and the philosophy of science. Despite the breadth of its foundations, this volume is designed to be self-contained. Readers are expected to have a basic understanding of mathematics, while some familiarity with computer programming may be helpful. The book is intended for a broad technical audience, including mathematicians, computer scientists, engineers, and other scientifically inclined readers. Its mathematical level is suitable for advanced undergraduates and graduate students.

This first volume is dedicated to the mathematical foundations of the Theory of Nescience. It is organized into seven chapters and three appendices. Following an introductory chapter, six chapters develop the fundamental elements of the theory. The appendices provide the mathematical and philosophical background and a reference guide to the notation used throughout the book. Readers familiar with discrete mathematics, computability, Kolmogorov complexity, and multiobjective optimization may begin directly with the theoretical chapters, consulting the background material as needed.

\emph{Chapter \ref{chap:Introduction} Introduction} provides a gentle introduction to the Theory of Nescience, presenting an overview of its fundamental concepts and principal results. It introduces the distinction between entities, representations, and descriptions, explains the two fundamental sources of ignorance, miscoding and mismodel, and outlines the concepts of nescience, perfect knowledge, and unknown unknowns. The chapter is intended to be accessible without advanced mathematical knowledge and may be read independently as a general introduction to the theory.

\emph{Chapter \ref{chap:Entities} Entities} introduces the objects about which scientific knowledge is sought. It defines the notion of a research entity, distinguishes between different types of entities, and explains why the set of entities under investigation must be fixed relative to a particular domain. The chapter also introduces research areas as collections of entities that can be investigated jointly.

\emph{Chapter \ref{chap:Representations} Representations} examines how research entities can be represented by finite binary strings. It introduces the representation oracle, distinguishes between decodable, minimal, admissible, extended, and effective representations, and investigates the conditions under which entities are knowable or unknowable. The chapter also studies operations for augmenting, reducing, and combining representations, and extends the framework to research areas.

\emph{Chapter \ref{chap:Descriptions} Descriptions} studies how representations can be modeled by computable descriptions. It distinguishes between decodable, exact, minimal, extended, and effective descriptions, and introduces the concepts of hypotheses and mathematical models. The chapter also examines operations for combining, simplifying, and conditioning descriptions, establishing the formal basis for evaluating their adequacy and complexity.

\emph{Chapter \ref{chap:Miscoding} Miscoding} introduces a mathematical measure of the inadequacy of a representation with respect to an entity. It distinguishes between representational deficiency, arising from missing information, and representational surplus, arising from unnecessary or irrelevant information. These components are combined into a measure called miscoding. The chapter investigates its mathematical properties, relative formulations, and conditions under which miscoding can be reduced, extending the analysis to effective representations and research areas.

\emph{Chapter \ref{chap:Mismodel} Mismodel} introduces a mathematical measure of the inadequacy of a description with respect to a representation. It distinguishes between descriptive deficiency, arising from insufficient information, and descriptive surplus, arising from unnecessary complexity. These components are combined into a measure called mismodel. The chapter studies its mathematical properties, relative formulations, and conditions under which mismodel can be reduced. It also extends the analysis to conditional descriptions and research areas.

\emph{Chapter \ref{chap:Nescience} Nescience} brings together the preceding concepts into a unified mathematical framework for quantifying ignorance. It formulates the pursuit of scientific knowledge as a multiobjective optimization problem in which miscoding and mismodel are minimized simultaneously. The chapter introduces Pareto optimality and scalar measures of nescience, examines the conditions for perfect knowledge, and investigates how nescience may change as representations and descriptions are refined. It also extends the concept of nescience to research areas.

\emph{Appendix \ref{chap:Mathematical_Background} Mathematical Background} provides a concise introduction to the mathematical concepts required to understand the theory. It reviews the fundamentals of discrete mathematics, computability, complexity, information distance, randomness, and multiobjective optimization. The appendix establishes the notation and mathematical results used throughout the book, while directing readers to standard references for more detailed treatments.

\emph{Appendix \ref{chap:Philosophy} Philosophical Foundations} places the Theory of Nescience within the broader philosophical study of scientific knowledge. It reviews fundamental questions concerning the nature of science, research entities, scientific representation, discovery, explanation, and justification. The appendix also examines the limits of scientific understanding, providing a philosophical context for the mathematical framework developed in the preceding chapters.

\emph{Appendix \ref{chap:Notation} Notation} provides a reference guide to the mathematical symbols and notation used throughout the book.

These chapters establish the theoretical foundations of the Theory of Nescience. The development of practical applications in artificial intelligence, scientific modeling, and the systematic discovery of unknown research problems forms part of the broader research program beyond the scope of this volume.

%
% Section: Acknowledgments
%

\section*{Acknowledgments}

I would like to express my gratitude to everyone who has contributed their comments, ideas, and encouragement to the development of the Theory of Nescience. In particular, I am grateful to Antonio Fernández, Vincenzo Mancuso, and Paolo Casari, who believed in and supported this project from its earliest stages, when it was little more than a speculative idea. Others who have provided valuable feedback and contributions include Héctor Cordobés, Luis F. Chiroque, Agustín Santos, Marco Ajmone, Pablo Rojo, Manuel Cebrián, Andrés Ortega, Emilio Amaya, Mattis Choummanivong, Alexander Lynch, Andrés Carrillo, Simon Bihoreau, and Ricardo Hernando.

Some of the early computational work associated with the development of the theory received partial funding from the IMDEA Networks Institute, the European Union's Horizon 2020 research and innovation programme under grant agreement No. 732667 (RECAP), and Nokia Spain through the NetPredict project.

The ideas presented in this book build upon a long tradition of mathematical, computational, and philosophical thought. I draw inspiration from philosophers such as Plato, Occam, Popper, Feyerabend, and Wittgenstein; from polymaths such as Llull, Leibniz, and Newton; and from the mathematical foundations established by Turing, Church, Post, Shannon, Solomonoff, Chaitin, Kolmogorov, and many others. My contribution lies in connecting and reinterpreting some of these insights within a unified mathematical framework. The references provided throughout the book acknowledge the works that have most directly influenced its development.

I am also grateful to the open-source community, whose contributions have made this work possible. The preparation of this book relied extensively on freely available tools, from \TeX{}, \LaTeX{}, and \texttt{TikZ} for writing and mathematical illustrations to Python and scientific libraries such as \texttt{NumPy}, \texttt{pandas}, \texttt{Matplotlib}, and \texttt{scikit-learn} for computational experiments. The dedication of countless developers to creating and maintaining these tools exemplifies the collaborative spirit upon which scientific progress depends.

I would also like to acknowledge the use of ChatGPT in the preparation of this book. As an AI assistant, it provided support in reviewing and refining the manuscript, improving clarity and language, suggesting examples, exploring alternative formulations, and assisting with the development and examination of mathematical arguments and proofs. All AI-generated suggestions were subject to my review, and I retain full responsibility for the mathematical content, conclusions, and any remaining errors. The opportunity to explore ideas through interactive dialogue has been a valuable part of the writing and research process.

Finally, I am deeply grateful to my parents for providing me with opportunities they themselves never had, and to my wife and three children, who give my life its deepest meaning.